\section{二阶结构区分局部极小与鞍点}
\label{sec:09-Convex-Optimization/08-Nonconvex-and-Stochastic-Optimization/02}

训练到第八千步，损失曲线横住了，最近两千步只在第四位小数上抖，梯度范数稳定在 \(3\times10^{-4}\)。会议室里有两种意见：一种说训练完了，收工；另一种说卡在鞍点上，把学习率调大踢一脚。

这两种意见对应的对策相反，代价也相反。判错第一种，你白烧几万卡时；判错第二种，你把一个已经不错的解踢出了它的盆地。而分歧的来源很具体——梯度范数在两种情形下读数一样。上一篇已经把这个局限交代清楚了：一阶信息只认证脚下平坦。要往下分，得再要一阶导数。

\subsection{驻点处的领头项是 Hessian 的二次型}

在驻点 \(x^\star\) 处展开，一阶项消失，剩下的领头项是二次型：

\[
f(x^\star+h)=f(x^\star)+\tfrac12 h^{\trans}\nabla^2 f(x^\star)h+o(\norm h^2).
\]

Hessian 是对称矩阵，可以正交对角化。特征向量给出主曲率方向，特征值给出沿该方向的弯曲程度和方向：正的向上弯，负的向下弯。沿第 \(i\) 个单位特征向量 \(v_i\) 走 \(t\) 步，函数值的变化是 \(\tfrac12\lambda_i t^2+o(t^2)\)。特征值的符号，就此决定这个驻点属于哪一类。

\begin{center}
\begin{tikzpicture}[x=1.1cm,y=1.1cm,font=\small,>=stealth]
  \foreach \a/\b in {0.42/0.30, 0.82/0.58, 1.22/0.86}{
    \draw[black!45] (0,0) ellipse [x radius=\a, y radius=\b];}
  \fill[black!80] (0,0) circle (1.6pt);
  \draw[->,black!70] (1.22,0) -- (0.55,0);
  \draw[->,black!70] (-1.22,0) -- (-0.55,0);
  \draw[->,black!70] (0,0.86) -- (0,0.42);
  \draw[->,black!70] (0,-0.86) -- (0,-0.42);
  \node[font=\footnotesize,black!70] at (0,-1.35) {局部极小};
  \node[font=\footnotesize,black!60] at (0,-1.72) {特征值全为正};
  \begin{scope}[shift={(3.0,0)}]
    \foreach \a/\b in {0.42/0.30, 0.82/0.58, 1.22/0.86}{
      \draw[black!45] (0,0) ellipse [x radius=\a, y radius=\b];}
    \fill[black!80] (0,0) circle (1.6pt);
    \draw[->,black!70] (0.55,0) -- (1.22,0);
    \draw[->,black!70] (-0.55,0) -- (-1.22,0);
    \draw[->,black!70] (0,0.42) -- (0,0.86);
    \draw[->,black!70] (0,-0.42) -- (0,-0.86);
    \node[font=\footnotesize,black!70] at (0,-1.35) {局部极大};
    \node[font=\footnotesize,black!60] at (0,-1.72) {特征值全为负};
  \end{scope}
  \begin{scope}[shift={(6.0,0)}]
    \foreach \sx in {1,-1}{\foreach \sy in {1,-1}{
      \begin{scope}[xscale=\sx,yscale=\sy]
        \draw[black!45] plot[domain=0.17:1.2,samples=50] (\x,{0.2/\x});
        \draw[black!45] plot[domain=0.42:1.2,samples=50] (\x,{0.5/\x});
        \draw[black!45] plot[domain=0.83:1.2,samples=40] (\x,{1.0/\x});
      \end{scope}}}
    \draw[black!25,dashed] (-1.2,-1.2) -- (1.2,1.2);
    \draw[black!25,dashed] (-1.2,1.2) -- (1.2,-1.2);
    \fill[black!80] (0,0) circle (1.6pt);
    \draw[->,black!70] (0.95,0.95) -- (0.42,0.42);
    \draw[->,black!70] (-0.95,-0.95) -- (-0.42,-0.42);
    \draw[->,black!70] (0.42,-0.42) -- (0.95,-0.95);
    \draw[->,black!70] (-0.42,0.42) -- (-0.95,0.95);
    \node[font=\footnotesize,black!70] at (0,-1.35) {鞍点};
    \node[font=\footnotesize,black!60] at (0,-1.72) {符号有正有负};
  \end{scope}
\end{tikzpicture}

\emph{等高线本身分不出前两种，箭头（下降方向）才分得出；第三种连等高线的形状都不同——闭合曲线换成了相交的双曲线族。}
\end{center}

图上前两幅的等高线一模一样，区别只在箭头指向中心还是背离中心。这不是画图偷懒：在极小点和极大点附近，等高面都是同心的闭曲面，拓扑上没有差别，差别写在函数值的高低次序里。第三幅不一样——等高线在临界点处交叉，把邻域切成四块，两块高两块低。鞍点的特征就在这个交叉上：存在一条路能下去，只是它不是任意一条。

\subsection{一个负特征值就足以排除局部极小}

必要条件是两条：\(\nabla f(x^\star)=0\) 且 \(\nabla^2 f(x^\star)\succeq0\)。充分条件把第二条加强为 \(\nabla^2 f(x^\star)\succ0\)。

\begin{proposition}
若 \(\nabla f(x^\star)=0\) 且 \(\nabla^2 f(x^\star)\) 有一个负特征值，则 \(x^\star\) 不是局部极小点。
\end{proposition}

骨架三步。取该负特征值对应的单位特征向量 \(v\)，则 \(v^{\trans}\nabla^2 f(x^\star)v=\lambda<0\)。沿 \(v\) 走 \(t\) 步，Taylor 展开给出 \(f(x^\star+tv)=f(x^\star)+\tfrac12\lambda t^2+o(t^2)\)。当 \(t\) 足够小时余项的绝对值小于 \(\tfrac14\abs{\lambda}t^2\)，于是 \(f(x^\star+tv)<f(x^\star)\)，任意小的邻域里都存在更低的点。

条件对账：用到了二阶连续可微（否则 Taylor 余项没有 \(o(\norm h^2)\) 的控制），用到了 \(\lambda<0\) 严格小于零（等于零时余项可能反号，\(f(x)=x^4\) 与 \(f(x)=-x^4\) 在原点的二阶行为完全相同而结论相反），没有用到凸性、没有用到 Hessian 的其他特征值。

这个判据在实践中比正定性判据好用，因为它是排除性的：找到一个负曲率方向就够了，不必检查所有方向。而半正定这个必要条件反过来什么也证明不了——\(f(x)=x^3\) 和 \(f(x)=x^4\) 在原点的一阶二阶行为完全一致，一个是拐点一个是严格极小，差别藏在三阶及以上。Hessian 退化的时候，二阶条件就交出了判断权。

\subsection{高维参数空间里全正是稀有事件}

现在把维数放进来。一个 \(d\) 维问题的 Hessian 有 \(d\) 个实特征值，局部极小要求它们全为正。粗暴地把每个符号当成独立的公平硬币，全正的概率是 \(2^{-d}\)。这个算法当然不严谨——特征值之间有强相关，符号也未必等概率——但它给出的量级方向是对的。

严谨一点的说法来自随机矩阵理论。把 Hessian 建模成 Wigner 型随机对称矩阵，谱在 \(d\) 大时集中到半圆分布上（见\hyperref[sec:11-Mathematical-Statistics/09-Random-Matrix-Theory/02]{``Marchenko--Pastur 律给出谱的精确形状''}），负特征值的个数集中在 \(d/2\) 附近，全部落到正半轴的概率随 \(d\) 指数衰减。经验风险曲面上还有一条更细的观察：这类模型里，函数值越低的临界点，负特征值比例越小；高处的临界点几乎全是鞍点，低处才开始出现真正的极小。

对一个参数量 \(10^8\) 的网络，这些说法合起来的操作性结论只有一条：当训练卡住时，先验上更可能是停在鞍点或宽平台附近，而不是撞上了一个糟糕的局部极小。上世纪关于神经网络``到处是坏局部极小''的担忧，很大一部分在维数变高之后被这个观察改写了。要说明的是，这些是在特定随机模型下算出来的启发式结论，不是对任意损失曲面的定理。

\subsection{逃离靠负曲率方向，或者靠噪声}

排除了局部极小之后，怎么走。

一条路是显式地找负曲率方向。参数上百万时 Hessian 存不下，但 Hessian 与向量的乘积可以由自动微分算出来——对 \(\ip{\nabla f(x)}{u}\) 再求一次梯度即得 \(\nabla^2f(x)u\)，代价约等于一次额外的反向传播。有了这个算子，Lanczos 迭代或在 \(-\nabla^2f\) 上跑幂法就能估计最小特征值和对应方向。这条路和 \hyperref[sec:09-Convex-Optimization/06-Optimization-Algorithms/02]{``Newton 法用曲率校正方向''}的算子视角一致：需要的是矩阵的作用效果，不是矩阵本身。

另一条路更便宜，实践中也更常见：让随机梯度自己带来的噪声去踢。如果目标函数满足 strict-saddle 结构——每个非极小驻点都有一个严格负曲率方向——那么在适当的光滑性条件下，加了扰动的梯度下降能以高概率在多项式时间内绕开所有鞍点。相当一批矩阵分解和相位恢复问题被证明具备这个结构。

浮点计算里这些条件都得放宽。实用的停止判据是近似二阶驻点：

\[
\norm{\nabla f(x)}\le\varepsilon_g,\qquad \lambda_{\min}\bigl(\nabla^2f(x)\bigr)\ge-\varepsilon_H .
\]

第一式说没有明显的一阶下降方向，第二式说没有能一步吃到大收益的负曲率出口。两式一起，比单看梯度范数强得多，但仍然只是局部证据——别处可能有更深的谷，它们对此沉默。

还有几条来自实践的``但是''要说清楚。负特征值可能小到 \(-10^{-8}\)，理论上的下降方向在有限精度和有限步长下几乎没有推力，算法在附近黏很久，看上去和收敛没有区别。平台也可能很宽，曲率在一大片区域内都接近零，等走到负曲率明显的地方，预算已经花掉大半。反过来，噪声调得太大，它把你踢出鞍点的同时也可能把你踢出一个不错的盆地——噪声尺度需要与地形匹配，而地形是未知的。

\subsection{零特征值多半来自参数化的冗余}

深度网络的 Hessian 谱里总有一大堆接近零的特征值，第一次看到会以为哪里出了问题。多数时候它反映的是参数化本身的对称性：交换同一层里两个神经元的编号，网络实现的函数一字不变；ReLU 网络里把一层的权重放大 \(c\) 倍、下一层缩小 \(c\) 倍，输出也一字不变。这类操作在参数空间里画出连续的路径，沿路径函数值恒定，对应方向上的曲率自然是零。

于是极小解常常不是孤立的点，而是一整条曲线甚至更高维的流形。Hessian 奇异在这里不是病症，只是坐标选得冗余。

要分辨的是另一件事：一个零曲率方向究竟是无害的参数冗余（动了参数不动预测），还是一个训练数据管不住、测试时会明显改变输出的薄弱方向。二阶谱本身答不了这个问题，它只描述训练损失的局部形状。答案要到数据和表示那一侧去找，这条线索留到本单元最后一篇。

上面这些都在讨论驻点有多难缠。但实际训练里还有一类相反的经验：某些明显非凸的问题跑起来异常顺，损失一路降到底，从没在什么鞍点上耽搁过。那类地形凭什么好走，下一篇给出一个条件。
